% \subsubsection{Generation of a model and specifications for formal-verification}
% A custom program has to be written in Python, that parses the Xacro files and ROS2 Nodes controller to generate a NuSMV model:
%
  % - a Boolean occupancy matrix representing the neighbourhood on a field relevant to the current replication.
%
  % - replication program-counter plus local flags (gripper-occupied, block -held, etc.).
%
  % The liveness and safety properties has to be defined independent from the simulation, in a form close to the following:
  % \begin{itemize}
    % 
      % \item Liveness. $G (start \rightarrow F (Robot_1==ready \land Robot_2==ready))$ — once replication starts, eventually both robots are able to function independently.
%
      % \item Collision-free safety. $G \neg collision$ (or $AG \neg collision$ for CTL).
%
% \end{itemize}

\newpage
\section{Robot configuration model checking}

\subsection{Overview}

In this part we will check the model of the robot of the certain configuration 
for the reachability property introduced in 

Block types -> NuSMV block modules structure -> NuSMV template\\
Border conditions:
    - Base blocks on the end
    - initial block position
    - connect only from upper side

python script -> Generate NuSMV -> check reachability border


Criteria: reachability of the random distributed blocks, given base block position
Validation: 
optimal structure -> generate xacro, check assembly


\subsection{Searching for relevant robot configurations}
\subsubsection{Individual robot configuration checking}
In order to systematically evaluate which robot structures are capable 
of reaching a target coordinates, implemented an automated configuration checking 
procedure that integrates model generation, model checking with NuSMV, and 
result filtering. The overall workflow is illustrated in Figure 
\ref{fig:SearchingRelevantConfigurations}.

\begin{figure}[h!]
    \centering
    \includegraphics[height=15cm]{SearchingRelevantConfigurations.png} 
    \caption{Activity diagram of searching for relevant configurations}
    \label{fig:SearchingRelevantConfigurations}
\end{figure}

The process of checking robot configurations is performed by integrating Python
-based model generation with the NuSMV model checker. The workflow begins 
with two primary inputs: the set of available block types, such as 
base, yaw-rotating, and pitch-rotating blocks, and the maximum robot length expressed as 
the number of blocks. These inputs are used to generate the complete set 
of allowed configurations of length $N$, ensuring that every possible arrangement 
of building blocks within the specified limit is considered. 
Each configuration is represented as an ordered sequence 
of blocks, starting with the mandatory base block and followed by a 
selection of yaw and pitch blocks.
Once the set of configurations is created, the process iterates over 
each configuration individually. \\
The block sequence is extracted and combined with:
\begin{itemize}
\item coordinates of an initial block of a robot ($x_{orig}$,$y_{orig}$,$z_{orig}$)
\item assembly site coordinates  relative to a robot($x_{asm}$,$y_{asm}$,$z_{asm}$)
\item target block coordinates relative to a robot:
 [Block($x_{1}$,$y_{1}$,$z_{1}$), Block2($x_{2}$,$y_{2}$,$z_{2}$), ...]
\end{itemize}
This information forms the basis for generating 
a concrete NuSMV model through a Python script. The script instantiates the 
appropriate modules corresponding to each block type, establishes the 
connections between their input and output variables, and integrates the 
positional and orientation parameters across the robot chain. In addition
, the script embeds the temporal logic specifications provided as LTL 
or CTL formulas. These specifications define the verification goals, 
such as whether the end-effector can reach a given target coordinates and 
avoids certain unsafe areas.

The generated model is then passed to NuSMV, which executes symbolic model 
checking to evaluate whether the specifications are satisfied. The model 
checker systematically explores the entire state space of the configuration, 
taking into account the discrete transitions of yaw and pitch blocks as 
well as the coordinate propagation along the chain. \\ 
If the specification is not satisfied, the algorithm discards the configuration and proceeds to 
the next candidate in the sequence. \\ 
If the specification is satisfied, 
the configuration is stored in a set of relevant configurations, which 
represents the final output of the process.


\subsubsection{Implementation in Python}
The implementation of the  workflow described in 
Figure \ref{fig:SearchingRelevantConfigurations} is presented in three-class Python design.
The architecture of the design corresponds directly to the class diagram shown on the 
Figure \ref{fig:RobotConfigCheckingClassDiagram}, where the 
class \textit{Configuration\_checking} orchestrates the process, the class \textit{Robot} 
assembles a sequence of blocks, and the class \textit{Block} encapsulates the 
description of an individual block.

\begin{figure}[h!]
    \centering
    \includegraphics[height=10cm]{RobotConfigCheckingClassDiagram.png} 
    \caption{Class diagram of searching for relevant configurations}
    \label{fig:RobotConfigCheckingClassDiagram}
\end{figure}
A \textit{Block} class is the minimal self-contained representation of a kinematic segment 
in the robot block sequence. Each block is identified by an integer number and 
a  type such as “base”, “yaw”, or “pitch”, and it stores the 
index of its parent block in order to define the topology and orientation.
Geometric information is reduced to a single parameter length, while all 
pose and orientation information is described later through the NuSMV model.
The block provides three critical outputs in the form of NuSMV code 
fragments: a structure line that instantiates the variable with appropriate block module and 
wires it to its parent, an initialisation line that sets link length and 
default joint angle values, and, for the final block, assignment lines 
that expose the end-effector coordinates. For non-base blocks, angle variable 
names are generated systematically (e.g., yaw2, pitch3), guaranteeing uniqueness 
and readability. In this way, each block fully describes its structural role 
without knowledge of the entire robot, making the class simple and robust.

The Robot class acts as a container and assembler for blocks. Constructed from 
a list of block types, it automatically creates a chain of Block instances
, numbers them, and links each to its parent. It maintains the list of 
blocks as well as a handle to the final block, which represents the 
robot’s end-effector. The Robot then concatenates the fragments provided 
by its blocks into complete sections suitable for insertion into a 
NuSMV model. It can return the structure of the chain as a sequence of 
module instantiations, the initialisation of link lengths and angles, 
and the assignment of global end-effector coordinates. Additionally, 
it encodes the configuration in a readable name such as “base\_yaw\_pitch\_yaw”, 
which is used both in logging and file management. The Robot therefore translates 
an abstract sequence of types into a structured collection of blocks, 
and from there into concrete model code that describes the robot in full.

The third class, Configuration\_checking, governs the entire process. 
Its role is to generate and test all possible configurations of a 
specified chain length composed of a given set of block types. The first 
block is always fixed to a base, and the remaining positions are filled 
with every possible combination of the admissible types. This results in 
an exponential search space, for example two types and four positions yielding 
sixteen unique robots. Configuration\_checking takes as additional input 
a target point in space and stores both this point and the error tolerance 
around it, against which reachability is later checked. For each generated 
morphology, the class constructs a Robot, retrieves its NuSMV code, 
and splices it into a template file. The method create\_nusmv\_from\_template searches 
for the region between “MODULE main” and a predefined end marker, 
replaces that region with a new main module generated for the current robot
, and leaves the rest of the file untouched. This ensures that all library 
code, specifications, and fairness constraints in the template remain intact
. The new module contains the structure lines from the robot, the initialisation 
of lengths and angles, the fixed base pose and basis vectors, the end
-effector definitions, and the error-band predicates that compare the 
end coordinates to the desired target.
 This layering is explicit in the class diagram Figure\ref{fig:RobotConfigCheckingClassDiagram}: 
Configuration\_checking depends on Robot, which is composed of Block instances.

Instead of hard-coding a complete NuSMV 
model, the system rewrites  the MAIN module fragment inside a 
reusable template file. This makes the pipeline resilient to changes in 
property specifications or library modules—the template remains the 
source of truth, while the classes inject only the structure and constants.


The overall workflow combines automatic model generation, model checking,
and storing of satisfiable robot configurations. 
By using a Python script to handle tasks 
of configuration expansion and NuSMV code generation, the approach ensures 
scalability and reduces the potential errors. The inclusion of 
temporal logic specifications provides a basis for determining whether 
a configuration is functionally capable of achieving the task requirements. 
The final set of relevant configurations captures only those robot 
designs that satisfy the given conditions, providing a reliable foundation 
for further analysis or deployment in robotic applications.




\subsubsection{Aggregated robot configuration checking in NuSMV}

\newpage
